\section{Data Elements}

\subsection{Pair}
Eg: \texttt{pair<int, int> x = {1,2};}\\
\begin{tabularx}{\linewidth}{lX}
\texttt{x.first} & Returns 1. \\
\texttt{x.second} & Returns 2.
\end{tabularx}

\subsection{Tuple}
Fixed length collection of heterogenous values.\\
Eg: \texttt{tuple<int, string, double> t;}\\

\textbf{Operations related to tuple}\\
\begin{tabularx}{\linewidth}{lX}
\texttt{get<i>} & Return $i^th$ value of t. Eg: \texttt{get<0>(t);}.\\
\texttt{make\_tuple()} & Returns tuple out of individual values.\\
& Eg: \texttt{make\_tuple(5, "Paul", 23.4)}.\\
\texttt{tie()} & Unpacking tuple.\\
& Eg: \\
& \texttt{tie(n, s, x) = t;}\\
& \texttt{tie(n, std::ignore, x) = t;}\\
\end{tabularx}

\subsection{Bit set}

\textbf{Definiton}\\
\texttt{bitset<size> s;}\\
Examples:\\
\texttt{bitset<5> s;}
\texttt{bitset<5> s(string("1100110"));}\\
\textbf{Operations on Bit set:}\\
\begin{tabularx}{\linewidth}{lX}
\texttt{a \& b} & Bitwise AND.\\
\texttt{a | b} & Bitwise OR.\\
\texttt{a \^ b} & Bitwise EXOR.\\
\texttt{~a} & Bitwise negation.\\
\end{tabularx}
